% 319_Proton_Torus_En_ch.tex
% The Proton as a Vibrating Torus:
% Geometric Foundations of Mass Emergence in FFGFT
% Doc. 319, August 2026

\section*{Summary}

In FFGFT the proton is not described as a collection of quarks and
gluons in a volume, but as a \emph{vibrating torus} $T^4/\mathbb{Z}_3$.
Quarks and gluons are different mode types of this structure. The proton
mass is the vibrational energy of the torus. This document develops the
geometric picture, anchors it in the corpus, and marks all statements
by the epistemic standard [K]/[S]/[SETZUNG].

\begin{tcolorbox}[colback=blue!5!white,colframe=blue!60!black,
  title=Reading note: inner perspective]
\textbf{[K]} We describe the proton torus \emph{from inside} --
as part of the same $T^4$ system (Doc.~248). There is no outside
perspective. What ``confinement'' or ``freedom'' means is relational:
it describes the relation between the observer and the respective
topological system, not an intrinsic property of the observed field.
\end{tcolorbox}

\section{The Geometric Picture}

\textbf{[K]} FFGFT starts from a compact 4D torus $T^4$ with
$\mathbb{Z}_3$ orbifold symmetry. The fundamental relation is
(Doc.~003):
\begin{equation}
  \tilde{T} \cdot m = 1
\end{equation}
Intrinsic time and mass are dual. This has immediate consequences for
all modes on the torus.

\textbf{[K]} The triality of the D4 lattice is built into the lattice
geometry, not postulated (Doc.~314): the order-3 elements in
$\text{Aut}(D4)$ are \emph{half-integer} -- the glue-vector mixture
$(\tfrac{1}{2},\tfrac{1}{2},\tfrac{1}{2},\tfrac{1}{2})$ enters. The
$T^4/\mathbb{Z}_3$ structure with exactly 9 fixed points follows from
the D4 geometry. These half-integer elements are the geometric origin
of the fermionic statistics of the quark modes.

\section{Two Mode Types: Quarks and Gluons}

\textbf{[K]} The compact $\mathbb{Z}_3$ fibre carries the discrete,
mass-bearing spectrum -- not the unrolled axis (Doc.~285):
\begin{quote}
``The discrete, mass-bearing spectrum sits on the compact
$\mathbb{Z}_3$ fibre, not on the unrolled axis.''
\end{quote}

\begin{table}[h]
\centering
\caption{Quark and gluon modes on the torus [K]}
\begin{tabular}{lllll}
\toprule
Mode & Spectrum & Mass & Proper time $\tilde{T}$ & Doc. \\
\midrule
Quark modes & discrete, $\mathbb{Z}_3$ fibre
  & $m > 0$ & $\tilde{T} < \infty$ & 006, 285, 314 \\
Gluon modes & continuous, locally flat
  & $m = 0$ & $\tilde{T} = \infty$ & 003, 270 \\
\bottomrule
\end{tabular}
\end{table}

\textbf{[K]} Gluons are massless ($m = 0$), hence $\tilde{T} = \infty$
-- they perform no measurement act (Doc.~270). Locally, flatness holds:
``in the spatial limit, local flatness'' (Doc.~270). This means: locally
the massless gluon sees $\mathbb{R}^3$, but that is a local
homeomorphism, not a global statement.

\section{The Relational Nature of Confinement and Freedom}

\textbf{[K]} Unrolling is a measurement act of the embedded observer
(Doc.~270): ``The measurement chooses one circle as time and unrolls
it; a physical act, not a mathematical limit.''

\textbf{[K]} We (observers) are part of the $T^4$ system -- not outside
(Doc.~248): ``We are part of the system. We always describe from inside,
never from outside.''

This yields a fundamental insight for the photon-gluon distinction:

\begin{keyresult}[Confinement is relational [K/S]]
Photon and gluon are both massless ($m=0$), both timeless
($\tilde{T}=\infty$), both see $\mathbb{R}^3$ locally.
\textbf{Both carry boundary conditions.}

The difference is not intrinsic but \emph{relational}:
\begin{itemize}
  \item \textbf{Photon:} We (observers) are part of the same
    $T^4$ system. The photon inhabits the same topological sector.
    Its boundary condition is our boundary condition. It therefore
    appears to us as freely propagating.
  \item \textbf{Gluon:} We face the $\mathbb{Z}_3$ substructure
    of the proton torus from outside. The gluon inhabits a different
    topological sector. It appears to us as confined.
\end{itemize}
``Confinement'' is not a property of the gluon alone -- it is a
property of the relation observer $\leftrightarrow$ system
(Doc.~248, [K]). \textbf{[S]} The complete topological derivation of
why photon and gluon inhabit different sectors remains open.
\end{keyresult}

\section{Massless Field Modes: Locally Pure Energy}

\textbf{[K]} For massless fields ($m=0$), $\tilde{T}\cdot m=1$ is not
defined -- not as a limit, but exactly: on null geodesics $d\tau=0$;
along a light path no proper time elapses. There is no rest frame
(Doc.~312):
\begin{quote}
``Masslessness means $v=c$ exactly -- the saturation of the
inequality $v\le c$. There is no $c^*(\lambda)$;
$dc/d\lambda\equiv 0$.''
\end{quote}

\textbf{[K]} What remains locally is exclusively $E=pc$ -- pure energy
(Docs.~134, 312). This holds for photon field modes and gluon field
modes \emph{equally}. From a local point of view, every massless field
mode is pure energy. There is no ``gluon here'' or ``photon there'' in
an absolute sense -- only energy density in space.

\textbf{[K]} ``Particles'' are not fundamental in FFGFT (Docs.~049, 156):
``There are no particles at all! What we call particles are different
excitation patterns in the single field $\delta m(x,t)$.'' Photon and
gluon are field modes -- not objects.

\textbf{[S]} The difference photon/gluon lies not in the local
description (both: $E=pc$, both: $d\tau=0$, both: locally flat), but
in the \emph{global topological embedding} of the field mode:

\begin{center}
\begin{tabular}{lll}
\toprule
Field mode & Topology & Global behaviour \\
\midrule
Photon & $U(1)$ flux, extended & propagates in $T^4$ (our system) \\
Gluon & $\mathbb{Z}_3/SU(3)$ linking & standing wave in proton torus \\
\bottomrule
\end{tabular}
\end{center}

\textbf{[S]} The gluon field is not localisable: as a non-Abelian
$SU(3)$ field mode (Doc.~145: colour-anticolour linking of three flux
threads) it interacts with itself. There is no bounded ``gluon object''
-- only a global field configuration of the proton torus. The eight
gluon modes correspond to the eight non-trivial linking states of the
three flux threads = eight generators of $SU(3)_c$ (Doc.~145, [S]).

\section{Topological Phase Closure as Confinement Mechanism}

\textbf{[K]} On the closed torus a wave must complete an integer phase
shift over the circumference:
\begin{equation}
  \oint k \, \mathrm{d}l = 2\pi n, \quad n \in \mathbb{Z}
\end{equation}

\textbf{[K]} The $\mathbb{Z}_3$ action on $L^2(T^4)$ permutes the
mode orbits cyclically; the permutation matrix has eigenvalues
$\{1, \omega, \omega^2\}$, phases $0^\circ$, $+120^\circ$,
$-120^\circ$ -- structurally exactly the circulant diagonalisation of
the $\mathbb{C}^3$ fibre (Doc.~314). This enforces discrete modes
without walls.

\textbf{[S]} The step from this mode structure to full QCD confinement
dynamics has not yet been formally carried out.

\section{The Duality $\tilde{T} \cdot m = 1$ in the Proton}

\textbf{[K]} The complete energy relation
$E^2 = (m_0 c^2)^2 + (pc)^2$ (Doc.~318) shows: massless gluons carry
energy $E = pc$ despite $m_0 = 0$ -- they carry $\sim 99\,\%$ of the
proton mass as field energy.

\textbf{[K]} This is not a paradox. A standing wave has no proper time
($\tilde{T} = \infty$), but its energy $E = \hbar\omega$ is real and
measurable. The duality $\tilde{T} \cdot m = 1$ states: timelessness
and masslessness are equivalent. The \emph{system energy} of the bound
state -- not the eigen-energy of individual modes -- is the proton mass.

\begin{table}[h]
\centering
\caption{Time-bearing vs. timeless oscillation [K]}
\begin{tabular}{lll}
\toprule
Property & Quark (time-bearing) & Gluon (timeless) \\
\midrule
Proper time & $\tilde{T} > 0$ & $\tilde{T} = \infty$ \\
Energy & Rest mass $+$ motion & Pure field energy $E = pc$ \\
Spectrum & Discrete ($\mathbb{Z}_3$ fibre) & Locally continuous \\
Statistics & Fermion (D4 triality [K]) & Boson \\
\bottomrule
\end{tabular}
\end{table}

\section{Spin Statistics from D4 Triality}

\textbf{[K]} The fermionic statistics of the quark modes follows from
the D4 lattice geometry (Doc.~314): the 80 order-3 elements in
$\text{Aut}(D4)$ fall into three classes; the class with trace $-2$ is
\emph{half-integer} and fixed-point-free in $\mathbb{R}^4$. This
half-integrality is the geometric foundation of fermion statistics.

\textbf{[S]} The draft approach ``$(n_\theta + n_\phi)$ odd $\to$
fermion'' captures the result but not the path. Correct is: D4 triality
delivers half-integer representations, independently of the parity of
the winding numbers.

\section{Colour Charge and $\mathbb{Z}_3$ Structure}

\textbf{[K]} The $\mathbb{Z}_3$ triality of the D4 lattice is built
in, not postulated (Doc.~314). $\mathbb{Z}_3$ is the centre of $SU(3)$.

\textbf{[S]} Whether and how the full $SU(3)_c$ gauge structure emerges
from this has not yet been formally derived in the corpus.

\begin{tcolorbox}[colback=orange!8!white,colframe=orange!60!black,
  title=Status of the colour charge thesis [S]]
The statement ``colour charge is effective, not fundamental'' is a
geometric picture \textbf{[S]}, not a proved derivation. The
$\mathbb{Z}_3$ connection [K] is a structural hint at the $SU(3)_c$
structure -- the complete geometric derivation remains open.
\end{tcolorbox}

\section{$\Lambda_{\text{QCD}}$ and Minimal Frequency}

\textbf{[S]} The minimal frequency on a torus with radius $R \approx 1$~fm:
\begin{equation}
  \omega_{\min} = \frac{1}{R} \approx 197~\text{MeV}
  \approx \Lambda_{\text{QCD}}
\end{equation}
($1~\text{fm}^{-1} = 197.3$~MeV in natural units). Consistent with
Doc.~041 ($\Lambda_{\text{QCD}} = v\cdot\xi^{1/3} \approx 200$~MeV),
but not a complete derivation.

\section{The Proton as a Whole}

\textbf{[K]} The proton is not a container with particles but a
\emph{system quantity}: the quark modes (discrete, $\mathbb{Z}_3$ fibre,
massive) define the boundary condition. The gluon modes (locally
continuous, massless) fill the system with field energy. The proton
mass is the total energy of this coupled system.

\textbf{[S]} The geometric picture ``quarks at the boundary, gluons in
the volume'' is conceptually robust but not yet a formal derivation from
the torus geometry.

\section{Relation to the Corpus}

\begin{itemize}
  \item \textbf{Doc.~003}: $\tilde{T}\cdot m=1$, $D_\text{frak}=2.94$,
    $K_\text{frak}$ [K]
  \item \textbf{Doc.~248}: Embedded observer, inner perspective [K]
  \item \textbf{Doc.~270}: Measurement as unrolling, local flatness [K]
  \item \textbf{Doc.~285}: Discrete spectrum on $\mathbb{Z}_3$ fibre [K]
  \item \textbf{Doc.~311}: Rolled-up/unrolled, three ways [K]
  \item \textbf{Doc.~314}: D4 triality, half-integer,
    $T^4/\mathbb{Z}_3$ from D4 [K]
  \item \textbf{Doc.~318}: Proton as open bridge,
    $E^2=(m_0c^2)^2+(pc)^2$, scope [K]
  \item \textbf{Doc.~041}: $\Lambda_\text{QCD}\approx 200$~MeV [S]
  \item \textbf{Doc.~049}: Confinement formulation [S]
\end{itemize}

\section{Open Questions [S]}

\begin{enumerate}
  \item \textbf{$SU(3)_c$ emergence from $\mathbb{Z}_3$ triality.}
    The structural hint [K] is established. The full gauge structure
    derivation is missing.
  \item \textbf{Why photon and gluon inhabit different sectors.}
    The relational explanation (Doc.~248) is [K]. The formal
    topological derivation remains [S].
  \item \textbf{$\Lambda_\text{QCD}$ from torus geometry.}
    Doc.~041: order of magnitude [S]. Complete derivation open.
  \item \textbf{RGE bridge $\alpha_s(m_\tau)\to\alpha_s(M_Z)$.}
    Doc.~160: $\alpha_s(m_\tau)=3\xi^{1/4}$ [K]. Bridge open [S].
  \item \textbf{Formal boundary conditions on $T^4/\mathbb{Z}_3$.}
    Exact spectral theory still to be developed [S].
\end{enumerate}
